\import{set/equinumerosity.tex}

\section{Equinumerosity and Cartesian products}

\begin{theorem}\label{times_equinum_congruent}
    Suppose $A\approx C$ and $B\approx D$.
    Then $A\times B\approx C\times D$.
\end{theorem}
\begin{proof}
    Take $f$ such that $f\in\bijections{A}{C}$ by \obv{equinum}.
    Take $g$ such that $g\in\bijections{B}{D}$ by \obv{equinum}.
    $f\in\funs{A}{C}$ and $g\in\funs{B}{D}$
        by \obv{bijections_to_funs}.
    Hence $f$ and $g$ are functions and $\dom{f}=A$ and $\dom{g}=B$
        by \obv{funs_is_function,funs}.
    $f$ and $g$ are injective by \obv{bijections}.
    Let $h=\{(p,(f(\fst{p}),g(\snd{p})))\mid p\in A\times B\}$.
    $h$ is a relation by \auto{relation}.
    $h$ is right-unique by \auto{rightunique,pair_eq_iff}.
    $\dom{h}=A\times B$ by \auto{setext,dom_iff,pair_eq_iff}.
    Show that for every $p\in A\times B$ we have
        $h(p)\in C\times D$.
    \begin{subproof}
        Fix $p\in A\times B$.
        $\fst{p}\in A$ and $\snd{p}\in B$
            by \obv{fst_type,snd_type}.
        $f(\fst{p})\in C$ and $g(\snd{p})\in D$
            by \obv{funs_type_apply}.
        $h(p)=(f(\fst{p}),g(\snd{p}))$
            by \obv{function_apply_intro}.
        Follows by \auto{times_tuple_intro}.
    \end{subproof}
    Hence $h\in\funs{A\times B}{C\times D}$ by \obv{funs_intro}.
    Show that $h$ is injective.
    \begin{subproof}
        Show that for all $p,q$ such that $p\in\dom{h}$ and
            $q\in\dom{h}$ and $h(p)=h(q)$ we have $p=q$.
        \begin{subproof}
            Fix $p,q$.
            Assume $p\in\dom{h}$ and $q\in\dom{h}$ and $h(p)=h(q)$.
            Hence $p,q\in A\times B$.
            $h(p)=(f(\fst{p}),g(\snd{p}))$ and
                $h(q)=(f(\fst{q}),g(\snd{q}))$
                by \auto{function_apply_intro}.
            Hence $f(\fst{p})=f(\fst{q})$ and
                $g(\snd{p})=g(\snd{q})$
                by \auto{pair_eq_iff}.
            $\fst{p},\fst{q}\in A$ and $\snd{p},\snd{q}\in B$
                by \obv{fst_type,snd_type}.
            Hence $\fst{p}=\fst{q}$ by \auto{injective_function}.
            Hence $\snd{p}=\snd{q}$ by \auto{injective_function}.
            Follows by
                \auto{times_elem_is_tuple,pair_eq_pair_of_proj,pair_eq_iff}.
        \end{subproof}
        Follows by \auto{injective_function}.
    \end{subproof}
    Show that for every $z\in C\times D$ there exists $p\in A\times B$
        such that $h(p)=z$.
    \begin{subproof}
        Fix $z\in C\times D$.
        Take $c,d$ such that $z=(c,d)$
            by \auto{times_elem_is_tuple}.
        Hence $c\in C$ and $d\in D$ by \obv{times_tuple_elim}.
        Take $a\in A$ such that $f(a)=c$ by \auto{surj,bijections}.
        Take $b\in B$ such that $g(b)=d$ by \auto{surj,bijections}.
        $(a,b)\in A\times B$ by \obv{times_tuple_intro}.
        $h((a,b))=(f(a),g(b))$
            by \auto{function_apply_intro,fst_eq,snd_eq}.
    \end{subproof}
    Hence $h\in\surj{A\times B}{C\times D}$ by \obv{surj}.
    Hence $h\in\bijections{A\times B}{C\times D}$
        by \obv{bijections}.
    Follows by \obv{equinum}.
\end{proof}

\begin{theorem}\label{times_injection_congruent}
    Let $f\in\inj{A}{C}$.
    Let $g\in\inj{B}{D}$.
    Then there exists $h\in\inj{A\times B}{C\times D}$.
\end{theorem}
\begin{proof}
    $f\in\funs{A}{C}$ and $g\in\funs{B}{D}$ by \obv{inj}.
    Hence $f$ and $g$ are functions and $\dom{f}=A$ and $\dom{g}=B$
        by \obv{funs_is_function,funs}.
    $f$ and $g$ are injective by \auto{inj,injective_function}.
    Let $h=\{(p,(f(\fst{p}),g(\snd{p})))\mid p\in A\times B\}$.
    $h$ is a relation by \auto{relation}.
    $h$ is right-unique by \auto{rightunique,pair_eq_iff}.
    $\dom{h}=A\times B$ by \auto{setext,dom_iff,pair_eq_iff}.
    Show that for every $p\in A\times B$ we have
        $h(p)\in C\times D$.
    \begin{subproof}
        Fix $p\in A\times B$.
        $\fst{p}\in A$ and $\snd{p}\in B$ by \obv{fst_type,snd_type}.
        $f(\fst{p})\in C$ and $g(\snd{p})\in D$
            by \obv{funs_type_apply}.
        $h(p)=(f(\fst{p}),g(\snd{p}))$
            by \obv{function_apply_intro}.
        Follows by \auto{times_tuple_intro}.
    \end{subproof}
    Hence $h\in\funs{A\times B}{C\times D}$ by \obv{funs_intro}.
    Show that for all $p,q$ such that $p\in\dom{h}$ and
        $q\in\dom{h}$ and $h(p)=h(q)$ we have $p=q$.
    \begin{subproof}
        Fix $p,q$.
        Assume $p\in\dom{h}$ and $q\in\dom{h}$ and $h(p)=h(q)$.
        Hence $p,q\in A\times B$.
        $h(p)=(f(\fst{p}),g(\snd{p}))$ and
            $h(q)=(f(\fst{q}),g(\snd{q}))$
            by \auto{function_apply_intro}.
        Hence $f(\fst{p})=f(\fst{q})$ and
            $g(\snd{p})=g(\snd{q})$ by \auto{pair_eq_iff}.
        $\fst{p},\fst{q}\in A$ and $\snd{p},\snd{q}\in B$
            by \obv{fst_type,snd_type}.
        Hence $\fst{p}=\fst{q}$ by \obv{injective_function}.
        Hence $\snd{p}=\snd{q}$ by \obv{injective_function}.
        Follows by
            \auto{times_elem_is_tuple,pair_eq_pair_of_proj,pair_eq_iff}.
    \end{subproof}
    Hence $h$ is injective by \auto{injective_function}.
    Hence $h\in\inj{A\times B}{C\times D}$
        by \obv{inj,injective_function}.
\end{proof}

\begin{proposition}\label{times_equinum_commutative}
    $A\times B\approx B\times A$.
\end{proposition}
\begin{proof}
    Let $s=\{(p,(\snd{p},\fst{p}))\mid p\in A\times B\}$.
    $s$ is a relation by \auto{relation}.
    $s$ is right-unique by \auto{rightunique,pair_eq_iff}.
    $\dom{s}=A\times B$ by \auto{setext,dom_iff,pair_eq_iff}.
    Show that for every $p\in A\times B$ we have $s(p)\in B\times A$.
    \begin{subproof}
        Fix $p\in A\times B$.
        $\fst{p}\in A$ and $\snd{p}\in B$ by \obv{fst_type,snd_type}.
        $s(p)=(\snd{p},\fst{p})$ by \obv{function_apply_intro}.
        Follows by \auto{times_tuple_intro}.
    \end{subproof}
    Hence $s\in\funs{A\times B}{B\times A}$ by \obv{funs_intro}.
    Show that $s$ is injective.
    \begin{subproof}
        Show that for all $p,q$ such that $p\in\dom{s}$ and
            $q\in\dom{s}$ and $s(p)=s(q)$ we have $p=q$.
        \begin{subproof}
            Fix $p,q$.
            Assume $p\in\dom{s}$ and $q\in\dom{s}$ and $s(p)=s(q)$.
            Hence $p,q\in A\times B$.
            $s(p)=(\snd{p},\fst{p})$ and
                $s(q)=(\snd{q},\fst{q})$
                by \auto{function_apply_intro}.
            Hence $\fst{p}=\fst{q}$ and $\snd{p}=\snd{q}$
                by \auto{pair_eq_iff}.
            Follows by
                \auto{times_elem_is_tuple,pair_eq_pair_of_proj,pair_eq_iff}.
        \end{subproof}
        Follows by \auto{injective_function}.
    \end{subproof}
    Show that for every $q\in B\times A$ there exists $p\in A\times B$
        such that $s(p)=q$.
    \begin{subproof}
        Fix $q\in B\times A$.
        Take $b,a$ such that $q=(b,a)$
            by \auto{times_elem_is_tuple}.
        Hence $a\in A$ and $b\in B$ by \auto{times_tuple_elim}.
        $(a,b)\in A\times B$ by \obv{times_tuple_intro}.
        $s((a,b))=(b,a)$
            by \obv{function_apply_intro,fst_eq,snd_eq}.
    \end{subproof}
    Hence $s\in\surj{A\times B}{B\times A}$ by \obv{surj}.
    Hence $s\in\bijections{A\times B}{B\times A}$
        by \obv{bijections}.
    Follows by \obv{equinum}.
\end{proof}

\begin{proposition}\label{times_equinum_associative}
    $(A\times B)\times C\approx A\times(B\times C)$.
\end{proposition}
\begin{proof}
    Let $r=\{(p,(\fst{\fst{p}},(\snd{\fst{p}},\snd{p})))\mid
        p\in(A\times B)\times C\}$.
    $r$ is a relation by \auto{relation}.
    $r$ is right-unique by \auto{rightunique,pair_eq_iff}.
    $\dom{r}=(A\times B)\times C$
        by \auto{setext,dom_iff,pair_eq_iff}.
    Show that for every $p\in(A\times B)\times C$ we have
        $r(p)\in A\times(B\times C)$.
    \begin{subproof}
        Fix $p\in(A\times B)\times C$.
        $\fst{p}\in A\times B$ and $\snd{p}\in C$
            by \obv{fst_type,snd_type}.
        Hence $\fst{\fst{p}}\in A$ and $\snd{\fst{p}}\in B$
            by \obv{fst_type,snd_type}.
        $r(p)=(\fst{\fst{p}},(\snd{\fst{p}},\snd{p}))$
            by \obv{function_apply_intro}.
        Follows by \auto{times_tuple_intro}.
    \end{subproof}
    Hence $r\in\funs{(A\times B)\times C}{A\times(B\times C)}$
        by \obv{funs_intro}.
    Show that $r$ is injective.
    \begin{subproof}
        Show that for all $p,q$ such that $p\in\dom{r}$ and
            $q\in\dom{r}$ and $r(p)=r(q)$ we have $p=q$.
        \begin{subproof}
            Fix $p,q$.
            Assume $p\in\dom{r}$ and $q\in\dom{r}$ and $r(p)=r(q)$.
            Hence $p,q\in(A\times B)\times C$.
            $r(p)=(\fst{\fst{p}},(\snd{\fst{p}},\snd{p}))$ and
                $r(q)=(\fst{\fst{q}},(\snd{\fst{q}},\snd{q}))$
                by \auto{function_apply_intro}.
            Hence $\fst{\fst{p}}=\fst{\fst{q}}$ and
                $\snd{\fst{p}}=\snd{\fst{q}}$ and
                $\snd{p}=\snd{q}$ by \auto{pair_eq_iff}.
            Hence $\fst{p}=\fst{q}$
                by
                \auto{fst_type,times_elem_is_tuple,pair_eq_pair_of_proj,pair_eq_iff}.
            Follows by
                \auto{times_elem_is_tuple,pair_eq_pair_of_proj,pair_eq_iff}.
        \end{subproof}
        Follows by \auto{injective_function}.
    \end{subproof}
    Show that for every $z\in A\times(B\times C)$ there exists
        $p\in(A\times B)\times C$ such that $r(p)=z$.
    \begin{subproof}
        Fix $z\in A\times(B\times C)$.
        Take $a,v$ such that $z=(a,v)$
            by \auto{times_elem_is_tuple}.
        Hence $a\in A$ and $v\in B\times C$ by \obv{times_tuple_elim}.
        Take $b,c$ such that $v=(b,c)$
            by \auto{times_elem_is_tuple}.
        Hence $b\in B$ and $c\in C$ by \obv{times_tuple_elim}.
        $((a,b),c)\in(A\times B)\times C$
            by \obv{times_tuple_intro}.
        $r(((a,b),c))=(a,(b,c))$
            by \auto{function_apply_intro,fst_eq,snd_eq}.
    \end{subproof}
    Hence $r\in\surj{(A\times B)\times C}{A\times(B\times C)}$
        by \obv{surj}.
    Hence $r\in\bijections{(A\times B)\times C}{A\times(B\times C)}$
        by \obv{bijections}.
    Follows by \obv{equinum}.
\end{proof}

\begin{proposition}\label{times_singleton_right_equinum}
    $A\times\{c\}\approx A$.
\end{proposition}
\begin{proof}
    Let $u=\{(p,\fst{p})\mid p\in A\times\{c\}\}$.
    $u$ is a relation by \auto{relation}.
    $u$ is right-unique by \auto{rightunique,pair_eq_iff}.
    $\dom{u}=A\times\{c\}$ by \auto{setext,dom_iff,pair_eq_iff}.
    For every $p\in A\times\{c\}$ we have $u(p)\in A$
        by \auto{function_apply_intro,fst_type}.
    Hence $u\in\funs{A\times\{c\}}{A}$ by \obv{funs_intro}.
    Show that $u$ is injective.
    \begin{subproof}
        Show that for all $p,q$ such that $p\in\dom{u}$ and
            $q\in\dom{u}$ and $u(p)=u(q)$ we have $p=q$.
        \begin{subproof}
            Fix $p,q$.
            Assume $p\in\dom{u}$ and $q\in\dom{u}$ and $u(p)=u(q)$.
            Hence $p,q\in A\times\{c\}$.
            $u(p)=\fst{p}$ and $u(q)=\fst{q}$
                by \auto{function_apply_intro}.
            $\snd{p}=c$ and $\snd{q}=c$
                by \auto{snd_type,singleton_iff}.
            Follows by
                \auto{times_elem_is_tuple,pair_eq_pair_of_proj,pair_eq_iff}.
        \end{subproof}
        Follows by \auto{injective_function}.
    \end{subproof}
    Show that for every $a\in A$ there exists $p\in A\times\{c\}$
        such that $u(p)=a$.
    \begin{subproof}
        Fix $a\in A$.
        $(a,c)\in A\times\{c\}$
            by \auto{times_tuple_intro,singleton_iff}.
        $u((a,c))=a$ by \obv{function_apply_intro,fst_eq}.
    \end{subproof}
    Hence $u\in\surj{A\times\{c\}}{A}$ by \obv{surj}.
    Hence $u\in\bijections{A\times\{c\}}{A}$
        by \obv{bijections}.
    Follows by \obv{equinum}.
\end{proof}

\begin{corollary}\label{times_singleton_left_equinum}
    $\{c\}\times A\approx A$.
\end{corollary}
\begin{proof}
    $\{c\}\times A\approx A\times\{c\}$
        by \obv{times_equinum_commutative}.
    $A\times\{c\}\approx A$ by \obv{times_singleton_right_equinum}.
    Follows by \obv{equinum_tran}.
\end{proof}

\begin{proposition}\label{singletons_equinum}
    $\{a\}\approx\{b\}$.
\end{proposition}
\begin{proof}
    $\{a\}\times\{b\}\approx\{a\}$
        by \obv{times_singleton_right_equinum}.
    $\{a\}\times\{b\}\approx\{b\}$
        by \obv{times_singleton_left_equinum}.
    Follows by \obv{equinum_sym,equinum_tran}.
\end{proof}
